\section*{Capítulo \ref{ch:b1:arith} --- Aritmética de los enteros}

\begin{solution}{exo:b1:arith:1}
$1001 = 1 \times 777 + 224$; $777 = 3 \times 224 + 105$;
$224 = 2 \times 105 + 14$; $105 = 7 \times 14 + 7$;
$14 = 2 \times 7 + 0$. Luego $\gcd(1001, 777) = 7$. Hacia atrás:
\[
7 = 105 - 7 \times 14
= 105 - 7(224 - 2\times 105) = 15 \times 105 - 7 \times 224
\]
\[
= 15(777 - 3\times 224) - 7\times 224 = 15 \times 777 - 52 \times 224
= 15 \times 777 - 52(1001 - 777) = 67 \times 777 - 52 \times 1001 .
\]
Comprobación: $67 \times 777 = 52\,059$ y $52 \times 1001 = 52\,052$;
la diferencia es $7$. Par de Bézout: $(u, v) = (-52, 67)$ para
$1001u + 777v = 7$.
\end{solution}

\begin{solution}{exo:b1:arith:2}
Como $10 \equiv 1 \pmod 9$: $10^k \equiv 1$, luego $\sum_k d_k 10^k
\equiv \sum_k d_k \pmod 9$. Como $10 \equiv -1 \pmod{11}$: $10^k
\equiv (-1)^k$, luego el entero es congruente con la suma alternada
$\sum_k (-1)^k d_k$ \omterm{def:b1:complex:field}{módulo} $11$ (empezando por la cifra de las
\emph{unidades} con signo $+$).

$123\,456\,789$: suma de cifras $45 \equiv 0 \pmod 9$. Suma alternada
desde las unidades: $9 - 8 + 7 - 6 + 5 - 4 + 3 - 2 + 1 = 5$, luego el
número es $\equiv 5 \pmod{11}$.
\end{solution}

\begin{solution}{exo:b1:arith:3}
Euclides: $237 = 2 \times 91 + 55$; $91 = 1 \times 55 + 36$;
$55 = 1 \times 36 + 19$; $36 = 1 \times 19 + 17$;
$19 = 1 \times 17 + 2$; $17 = 8 \times 2 + 1$. Hacia atrás:
\[
1 = 17 - 8\times 2 = 17 - 8(19 - 17) = 9\times 17 - 8\times 19
= 9(36 - 19) - 8\times 19 = 9\times 36 - 17\times 19
\]
\[
= 9\times 36 - 17(55 - 36) = 26\times 36 - 17\times 55
= 26(91 - 55) - 17\times 55 = 26\times 91 - 43\times 55
\]
\[
= 26\times 91 - 43(237 - 2\times 91) = 112 \times 91 - 43 \times 237.
\]
Así pues, $91 \times 112 \equiv 1 \pmod{237}$: las soluciones son
$x \equiv 112 \pmod{237}$. (Comprobación:
$91 \times 112 = 10\,192 = 43 \times 237 + 1$.)
\end{solution}

\begin{solution}{exo:b1:arith:4}
$\gcd(17, 39) = 1$: Euclides da $39 = 2\times 17 + 5$,
$17 = 3\times 5 + 2$, $5 = 2\times 2 + 1$, y hacia atrás
\[
1 = 5 - 2\times 2 = 5 - 2(17 - 3\times 5) = 7\times 5 - 2\times 17
= 7(39 - 2\times 17) - 2\times 17 = 7\times 39 - 16\times 17 .
\]
Solución particular $(x_0, y_0) = (-16, 7)$. Solución general de la
ecuación homogénea $17x + 39y = 0$: $x = 39k$, $y = -17k$ (pues
$17 \mid 39y$ y $\gcd(17,39) = 1$ obligan a $17 \mid y$ --- lema de
Gauss). Por tanto
\[
(x, y) = (-16 + 39k,\; 7 - 17k), \qquad k \in \Z .
\]
Para el miembro derecho $5$, multiplíquese por $5$ la solución
particular: $(x, y) = (-80 + 39k,\; 35 - 17k)$, $k \in \Z$.
\end{solution}

\begin{solution}{exo:b1:arith:5}
Para todo \omterm{def:b1:arith:prime}{primo} $p$, con $\alpha = v_p(a)$ y $\beta = v_p(b)$:
\[
v_p\bigl(\gcd(a,b)\bigr) + v_p\bigl(\operatorname{lcm}(a,b)\bigr)
= \min(\alpha, \beta) + \max(\alpha, \beta)
= \alpha + \beta = v_p(ab) .
\]
Dos enteros positivos con la misma valuación en todo \omterm{def:b1:arith:prime}{primo} son
iguales (\cref{prop:b1:arith:valuation}), luego
$\gcd(a,b)\operatorname{lcm}(a,b) = ab$.
\end{solution}

\begin{solution}{exo:b1:arith:6}
Valuaciones: $\gcd(a, b) = 2^{\min(10,6)} 3^{\min(4,7)} 5^{\min(2,0)}
7^{\min(0,1)} = 2^6\, 3^4 = 5184$;
$\operatorname{lcm}(a,b) = 2^{10}\, 3^7\, 5^2\, 7$.

Número de divisores: un divisor positivo de $n = \prod p_i^{\alpha_i}$
es exactamente una elección $\prod p_i^{\beta_i}$ con
$0 \leq \beta_i \leq \alpha_i$ (\cref{prop:b1:arith:valuation}); las
elecciones son independientes, luego hay $\prod_i (\alpha_i + 1)$
divisores. Para $a$: $(10+1)(4+1)(2+1) = 165$.
\end{solution}

\begin{solution}{exo:b1:arith:7}
Supóngase $\sqrt p = \frac rq$ con $r, q \in \N^*$, es decir,
$p q^2 = r^2$. Aplíquese $v_p$: $v_p(pq^2) = 1 + 2v_p(q)$ es impar,
mientras que $v_p(r^2) = 2 v_p(r)$ es par. Un entero no puede tener a
la vez valuación $p$-ádica par e impar: contradicción. Luego
$\sqrt p \notin \Q$.
\end{solution}

\begin{solution}{exo:b1:arith:8}
\emph{Hecho general.} Con $\gcd(m,n) = 1$, Bézout da $mu + nv = 1$.
Póngase $x_0 = b\,mu + a\,nv$. Entonces $x_0 \equiv a\,nv \equiv
a(1 - mu) \equiv a \pmod m$ y, análogamente, $x_0 \equiv b \pmod n$:
existencia. Si $x$ y $x'$ son dos soluciones, $m$ y $n$ \omterm{def:b1:arith:divides}{dividen} a
$x - x'$, luego $mn \mid x - x'$ (\cref{thm:b1:arith:gauss}~(2)):
unicidad \omterm{def:b1:complex:field}{módulo} $mn$.

\emph{Numéricamente:} $m = 7$, $n = 11$: $7 \times (-3) + 11 \times 2 =
1$. Luego $x_0 = 5 \times 7 \times (-3) + 2 \times 11 \times 2 = -105 +
44 = -61 \equiv 16 \pmod{77}$. Comprobación: $16 = 2\times 7 + 2
\equiv 2 \pmod 7$; $16 = 11 + 5 \equiv 5 \pmod{11}$. Soluciones:
$x \equiv 16 \pmod{77}$.
\end{solution}

\begin{solution}{exo:b1:arith:9}
\omterm{def:b1:complex:field}{Módulo} $7$: Fermat da $3^6 \equiv 1$, y $1000 = 6 \times 166 + 4$,
luego $3^{1000} \equiv 3^4 = 81 \equiv 4 \pmod 7$.

Dos últimas cifras de $7^{100}$: trabájese \omterm{def:b1:complex:field}{módulo} $4$ y \omterm{def:b1:complex:field}{módulo} $25$.
\omterm{def:b1:complex:field}{Módulo} $4$: $7 \equiv -1$, luego $7^{100} \equiv 1$. \omterm{def:b1:complex:field}{Módulo} $25$:
$7^2 = 49 \equiv -1$, luego $7^4 \equiv 1$ y
$7^{100} = (7^4)^{25} \equiv 1$. Por el teorema chino del resto
(\cref{exo:b1:arith:8}), $7^{100} \equiv 1 \pmod{100}$: las dos
últimas cifras son $01$.
\end{solution}

\begin{solution}{exo:b1:arith:10}
Escríbase $m = nq + r$, $0 \leq r < n$. Entonces
\[
2^m - 1 = 2^r\bigl(2^{nq} - 1\bigr) + 2^r - 1,
\]
y $2^n - 1$ \omterm{def:b1:arith:divides}{divide} a $2^{nq} - 1 = (2^n - 1)(2^{n(q-1)} + \dots +
1)$. Así pues, \omterm{def:b1:complex:field}{módulo} $2^n - 1$ se tiene $\;2^m - 1 \equiv 2^r - 1$ y,
como $0 \leq 2^r - 1 < 2^n - 1$, este \emph{es} el resto euclídeo.

Por tanto, el \omterm{met:b1:arith:euclid}{algoritmo de Euclides} sobre el par $(2^m - 1, 2^n - 1)$
reproduce, exponente a exponente, el algoritmo sobre $(m, n)$: cada
paso de división sustituye $(m, n)$ por $(n, r)$ arriba y
$(2^m - 1, 2^n - 1)$ por $(2^n - 1, 2^r - 1)$ abajo. Arriba el
algoritmo termina en $\gcd(m,n)$, luego abajo termina en
$2^{\gcd(m,n)} - 1$.
\end{solution}

\begin{solution}{exo:b1:arith:11}
Resuélvase primero $x^2 \equiv 1 \pmod p$: $p \mid (x-1)(x+1)$, luego,
por el lema de Euclides, $x \equiv 1$ o $x \equiv -1 \pmod p$.

En el producto $(p-1)! = 1 \times 2 \times \dots \times (p-1)$, todo
factor $a$ es invertible \omterm{def:b1:complex:field}{módulo} $p$, y su inverso $a^{-1}$ es de nuevo
uno de los factores (\cref{prop:b1:arith:invmod}). Emparéjese cada $a$
con $a^{-1}$: cada pareja multiplica a $1$, salvo los factores
emparejados consigo mismos ($a = a^{-1}$, es decir, $a^2 \equiv 1$),
que quedan solos --- y esos son exactamente $1$ y $p - 1$. Por tanto
\[
(p-1)! \equiv 1 \times (p - 1) \equiv -1 \pmod p .
\]
(Para $p = 2$: $1! = 1 \equiv -1 \pmod 2$; el argumento del
emparejamiento degenera, pero el resultado se mantiene.)

\emph{Recíproco.} Sea $n \geq 2$ compuesto, $n = ab$ con
$1 < a \leq b < n$. Si $a < b$, los dos aparecen como factores
distintos de $(n-1)!$, luego $n \mid (n-1)!$ y
$(n-1)! \equiv 0 \not\equiv -1$. Si $a = b$ (es decir, $n = a^2$):
para $a \geq 3$, tanto $a$ como $2a$ son $< n$, luego
$n = a^2 \mid a \times 2a \mid (n-1)!$, misma conclusión; y para
$n = 4$, $(n-1)! = 6 \equiv 2 \not\equiv -1 \pmod 4$.
\end{solution}

\begin{solution}{exo:b1:arith:12}
\begin{enumerate}
  \item Inducción. Para $n = 1$: $F_0 = 3 = F_1 - 2 = 5 - 2$.
        Suponiendo $F_0\cdots F_{n-1} = F_n - 2$:
        \[
        F_0 \cdots F_n = (F_n - 2)F_n
        = \bigl(2^{2^n} - 1\bigr)\bigl(2^{2^n} + 1\bigr)
        = 2^{2^{n+1}} - 1 = F_{n+1} - 2 .
        \]
  \item Sean $m < n$ y $d = \gcd(F_m, F_n)$. Por (1), $F_m$
        \omterm{def:b1:arith:divides}{divide} a $F_n - 2$, luego $d$ \omterm{def:b1:arith:divides}{divide} a la vez a $F_n$ y a
        $F_n - 2$, y por tanto \omterm{def:b1:arith:divides}{divide} a $2$. Pero todo número de
        Fermat es impar, luego $d = 1$.
  \item Cada $F_n \geq 3$ tiene un divisor \omterm{def:b1:arith:prime}{primo} $p_n$ (primer paso
        del~\cref{thm:b1:arith:euclidprimes}). Si $m \neq n$, entonces
        $p_m \neq p_n$, pues un \omterm{def:b1:arith:prime}{primo} común dividiría a
        $\gcd(F_m, F_n) = 1$. La \omterm{def:b1:logic:map}{aplicación} $n \mapsto p_n$ es, por
        tanto, \omterm{def:b1:logic:inj}{inyectiva} de $\N$ en los \omterm{def:b1:arith:prime}{primos}: hay infinitos
        \omterm{def:b1:arith:prime}{primos}.
\end{enumerate}
\end{solution}

\begin{solution}{pb:b1:arith:1}
\textbf{1.} $10! = 3\,628\,800$: dos ceros finales. Valuaciones factor
a factor: las potencias de $2$ vienen de $2, 4 = 2^2, 6, 8 = 2^3, 10$,
en total $v_2(10!) = 1 + 2 + 1 + 3 + 1 = 8$; las de $5$, de $5$ y
$10$: $v_5(10!) = 2$. Ceros finales $= \min(v_2, v_5) = 2$,
coherente.

\textbf{2.} Escríbase la división euclídea
$\lfloor x\rfloor = nq + r$, $0 \leq r \leq n - 1$. Entonces
$x = nq + r + \{x\}$ con $0 \leq r + \{x\} < n$, luego
$\lfloor x/n \rfloor = q = \bigl\lfloor \lfloor x \rfloor / n \bigr\rfloor$.

\textbf{3.} Sean $\alpha = v_p(a) \leq \beta = v_p(b)$
(intercámbiense si hace falta) y escríbanse $a = p^\alpha a'$,
$b = p^\beta b'$ con $p \nmid a', b'$. Entonces
$a + b = p^\alpha\bigl(a' + p^{\beta - \alpha}b'\bigr)$, luego
$v_p(a + b) \geq \alpha = \min$. Si $\alpha < \beta$, el paréntesis es
$a' + p^{\beta-\alpha}b' \equiv a' \not\equiv 0 \pmod p$: la valuación
vale exactamente $\alpha$.

\textbf{4.} Los múltiplos de $m$ en $\intint1n$ son
$m, 2m, \dots, qm$, donde $q$ es el mayor entero con $qm \leq n$, es
decir, $q = \lfloor n/m \rfloor$.

\textbf{5.} Por la factorización única,
$v_p(n!) = \sum_{j=1}^{n} v_p(j)$. Cuéntese de otro modo: cada $j$
aporta $v_p(j) = \#\{k \geq 1 : p^k \mid j\}$, luego
\[
v_p(n!) = \sum_{j=1}^n \#\{k : p^k \mid j\}
= \sum_{k\geq1} \#\{j \leq n : p^k \mid j\}
= \sum_{k\geq1} \Bigl\lfloor \frac n{p^k} \Bigr\rfloor
\]
por la pregunta 4 --- la fórmula de Legendre. La suma es finita: los
términos con $p^k > n$ se anulan.

\textbf{6.} $v_5(1000!) = 200 + 40 + 8 + 1 = 249$ (divisiones por
$5, 25, 125, 625$); $v_2(1000!) = 500 + 250 + 125 + 62 + 31 + 15 +
7 + 3 + 1 = 994$. Ceros finales de $1000!$: cada cero consume un $2$ y
un $5$, luego hay $\min(994, 249) = 249$.

\textbf{7.} Con $n = \sum_i a_ip^i$, la pregunta 2 da
$\lfloor n/p^k \rfloor = \sum_{i \geq k} a_ip^{i-k}$ (se trunca el
desarrollo en base $p$). Sumando en $k \geq 1$ e intercambiando las
dos sumas finitas:
\[
v_p(n!) = \sum_{i\geq1} a_i \sum_{k=1}^{i} p^{i-k}
= \sum_{i\geq0} a_i\,\frac{p^i - 1}{p - 1}
= \frac{n - s_p(n)}{p - 1} .
\]

\textbf{8.} Para $p = 2$: $v_2(n!) = n - s_2(n)$. Como todo $n \geq 1$
tiene $s_2(n) \geq 1$, siempre $v_2(n!) \leq n - 1 < n$:
$2^n \nmid n!$. Y $v_2(n!) = n - 1$ si y solo si $s_2(n) = 1$, si y
solo si $n$ es una potencia de $2$.

\textbf{9.} $n$ tiene $\lfloor \log_p n \rfloor + 1$ cifras en base
$p$, cada una a lo sumo $p - 1$, luego
$1 \leq s_p(n) \leq (p-1)\bigl(\log_p(n) + 1\bigr)$. Sustituyendo en
la pregunta 7:
\[
\frac n{p-1} - \log_p(n) - 1 \;\leq\; v_p(n!) \;<\; \frac n{p-1},
\]
y dividiendo por $n$: $\frac{v_p(n!)}n \to \frac1{p-1}$.

\textbf{10.} $Z(n) - Z(n-1) = v_5(n!/(n-1)!) = v_5(n)$: el número de
ceros finales salta $v_5(n)$ en cada múltiplo de $5$ y es constante
entre ellos. $Z(24) = \lfloor24/5\rfloor = 4$ y $Z(25) = 5 + 1 = 6$:
en $n = 25$ el recuento salta de $4$ directamente a $6$
($v_5(25) = 2$) y, como $Z$ es no decreciente con $Z \leq 4$ antes y
$Z \geq 6$ después, el valor $5$ no se alcanza nunca: ningún factorial
termina en exactamente cinco ceros.

\textbf{11.} Escríbase $x = \lfloor x\rfloor + \{x\}$:
$\lfloor x + y\rfloor = \lfloor x\rfloor + \lfloor y\rfloor +
\lfloor \{x\} + \{y\}\rfloor$, y $0 \leq \{x\} + \{y\} < 2$ hace que
la última parte entera valga $0$ o $1$. Después, aplicando Legendre
tres veces,
\[
v_p\binom{m+n}m = v_p\bigl((m{+}n)!\bigr) - v_p(m!) - v_p(n!)
= \sum_{k\geq1}\Bigl(
\Bigl\lfloor\frac{m+n}{p^k}\Bigr\rfloor
- \Bigl\lfloor\frac{m}{p^k}\Bigr\rfloor
- \Bigl\lfloor\frac{n}{p^k}\Bigr\rfloor\Bigr),
\]
una suma finita de $0$ y $1$ (aplíquese la primera afirmación a
$x = m/p^k$, $y = n/p^k$).

\textbf{12.} Fíjese $k \geq 1$ y escríbanse $m = p^km_1 + m_0$,
$n = p^kn_1 + n_0$ con $0 \leq m_0, n_0 < p^k$ (división euclídea:
$m_0$ es el número formado por las $k$ cifras bajas de $m$). Entonces
\[
\Bigl\lfloor\frac{m+n}{p^k}\Bigr\rfloor
- \Bigl\lfloor\frac m{p^k}\Bigr\rfloor
- \Bigl\lfloor\frac n{p^k}\Bigr\rfloor
= \Bigl\lfloor\frac{m_0 + n_0}{p^k}\Bigr\rfloor ,
\]
que vale $1$ si $m_0 + n_0 \geq p^k$ y $0$ en caso contrario. Pero
$m_0 + n_0 \geq p^k$ dice precisamente que sumar las $k$ cifras bajas
de $m$ y de $n$ desborda hacia la posición $k$ --- un acarreo hacia la
posición $k$ en el algoritmo escolar de la suma. Sumando en $k$:
$v_p\binom{m+n}m$ es el número de acarreos de la suma $m + n$ en base
$p$. (Kummer, 1852.)

\textbf{13.} Aplíquese Kummer con $m = j$, $n = p^k - j$, de suma
$p^k = (1\underbrace{0\cdots0}_{k})_p$. Sea $a = v_p(j)$: las cifras
en base $p$ de $j$ en las posiciones $0, \dots, a-1$ son $0$ y la
cifra en la posición $a$ no es nula. Las cifras de $p^k - j$ por
debajo de la posición $a$ también son $0$
($p^k - j = p^a(p^{k-a} - j/p^a)$). En la posición $a$, las dos cifras
no nulas deben sumar $p$ (cifra resultante $0$): un acarreo; y en cada
posición $a+1, \dots, k-1$, las cifras más el acarreo entrante suman
$p$ (de nuevo cifra resultante $0$): el acarreo se propaga. En total,
$k - a$ acarreos, luego $v_p\binom{p^k}j = k - v_p(j)$. Para $k = 1$:
$v_p\binom pj = 1$ para $0 < j < p$, la \omterm{def:b1:arith:divides}{divisibilidad} usada en
el~\cref{thm:b1:arith:fermat}.

\textbf{14.} Por la forma digital (pregunta 7), usando
$s_2(2n) = s_2(n)$ (se añade una cifra cero):
\[
v_2\binom{2n}n = \bigl(2n - s_2(2n)\bigr) - 2\bigl(n -
s_2(n)\bigr) = 2s_2(n) - s_2(2n) = s_2(n) \geq 1 :
\]
$\binom{2n}n$ es siempre par, y $v_2 = 1$ (es decir,
$\binom{2n}n \equiv 2 \pmod 4$) exactamente cuando $s_2(n) = 1$, es
decir, cuando $n$ es una potencia de $2$.

\textbf{15.} Vandermonde con $m = n = k = p$: $\binom{2p}p =
\sum_{j=0}^p \binom pj\binom p{p-j} = \sum_{j=0}^p \binom pj^2$. Para
$0 < j < p$, $p \mid \binom pj$ (pregunta 13), luego
$\binom pj^2 \equiv 0 \pmod p$; los términos de los extremos dan
$1 + 1$: $\binom{2p}p \equiv 2 \pmod p$.

\textbf{16.} Base $3$: $500 = 486 + 9 + 3 + 2$, cifras (de baja a
alta) $(2, 1, 1, 0, 0, 2)$, luego $s_3(500) = 6$; y
$1000 = 729 + 243 + 27 + 1$, cifras $(1, 0, 0, 1, 0, 1, 1)$, luego
$s_3(1000) = 4$. \emph{Kummer:} súmese $500 + 500$ en base $3$:
posición $0$: $2 + 2 = 4$, cifra $1$ y acarreo $1$; posición $1$:
$1 + 1 + 1 = 3$, cifra $0$ y acarreo $1$; posición $2$:
$1 + 1 + 1 = 3$, cifra $0$ y acarreo $1$; posición $3$:
$0 + 0 + 1 = 1$, sin acarreo; posición $4$: $0$; posición $5$:
$2 + 2 = 4$, cifra $1$ y acarreo $1$; posición $6$: cae el acarreo,
cifra $1$. Cuatro acarreos: $v_3\binom{1000}{500} = 4$.
\emph{Legendre:} $v_3(1000!) = \frac{1000 - 4}2 = 498$ y
$v_3(500!) = \frac{500 - 6}2 = 247$, luego
$v_3\binom{1000}{500} = 498 - 2\times247 = 4$. Los dos cálculos
coinciden --- y las cifras de la suma $(1, 0, 0, 1, 0, 1, 1)$
reproducen $1000$, como debe ser.

\textbf{17.} Por Kummer ($p = 2$, $m = k$, $n' = n - k$): $\binom nk$
es impar si y solo si la suma $k + (n - k)$ en base $2$ no tiene
acarreos, si y solo si en cada posición las cifras cumplen
$k_i + (n - k)_i = n_i$; en tal caso $k_i \leq n_i$ para todo $i$.
Recíprocamente, si $k_i \leq n_i$ para todo $i$, entonces el número de
cifras $n_i - k_i$ es $n - k$ y la suma no tiene acarreos. La misma
demostración en base $p$: $p \nmid \binom nk$ si y solo si cada cifra
en base $p$ de $k$ es a lo sumo la cifra correspondiente de $n$.

\textbf{18.} Contando los $k \in \intint0n$ cuyas cifras cumplen
$k_i \leq n_i$: cada cifra de $k$ se elige independientemente entre
$n_i + 1$ valores, lo que da $\prod_i (n_i + 1)$ elecciones; en base
$2$ esto es $2^{\#\{i : n_i = 1\}} = 2^{s_2(n)}$. Fila $4 = (100)_2$:
$2^1 = 2$ entradas impares --- en efecto, $1, 4, 6, 4, 1$ solo tiene
entradas impares en los extremos. Fila $5 = (101)_2$: $2^2 = 4$ ---
en efecto, $1, 5, 10, 10, 5, 1$.

\textbf{19.} Todas las entradas interiores son pares $\iff$ la fila
tiene exactamente $2$ entradas impares (las de los dos extremos
siempre lo son) $\iff 2^{s_2(n)} = 2 \iff s_2(n) = 1 \iff n$ es una
potencia de $2$.

\textbf{20.} En la suma de la pregunta 11, el $k$-ésimo término se
anula en cuanto $p^k > m + n$ (las tres partes enteras coinciden
entonces; de hecho la primera vale $0$ cuando $p^k > m+n$; más
simplemente, cada término es $0$). Por tanto, a lo sumo
$\lfloor \log_p(m+n)\rfloor$ términos son no nulos, y cada uno vale
$1$: $a = v_p\binom{m+n}m \leq \log_p(m+n)$, es decir,
$p^a \leq m + n$.

\textbf{21.} Para todo \omterm{def:b1:arith:prime}{primo} $p$,
$v_p\bigl(\operatorname{lcm}(1, \dots, 2n)\bigr) =
\lfloor\log_p(2n)\rfloor$ (la mayor potencia de $p$ que no supera a
$2n$ aparece entre $1, \dots, 2n$). La pregunta 20 con $m = n$ da
$v_p\binom{2n}n \leq \lfloor\log_p(2n)\rfloor$ para todo $p$: por
la~\cref{prop:b1:arith:valuation}, $\binom{2n}n \mid
\operatorname{lcm}(1, \dots, 2n)$. En cuanto al tamaño: el cociente
$\binom{2n}{k+1}/\binom{2n}k = \frac{2n-k}{k+1} \geq 1$ exactamente
para $k < n$, así que la entrada central es la mayor de las $2n + 1$
de la fila $2n$, de donde $4^n = \sum_k \binom{2n}k \leq
(2n+1)\binom{2n}n$. Combinando:
\[
\operatorname{lcm}(1, \dots, 2n) \geq \binom{2n}n \geq
\frac{4^n}{2n + 1} .
\]
Si hubiese pocos \omterm{def:b1:arith:prime}{primos} por debajo de $2n$, el mcm no podría ser tan
grande: el crecimiento exponencial del mcm es una huella cuantitativa
de la abundancia de \omterm{def:b1:arith:prime}{primos}.

\textbf{22.} $Z(n) = \sum_k\lfloor n/5^k\rfloor \approx \frac n4$,
así que se apunta cerca de $n = 4 \times 2026 = 8104$:
$Z(8104) = 1620 + 324 + 64 + 12 + 2 = 2022$. Súbase por múltiplos de $5$:
$Z(8110) = 2024$, $Z(8115) = 2025$, y
\[
Z(8120) = 1624 + 324 + 64 + 12 + 2 = 2026 .
\]
Como $Z$ es constante entre múltiplos de $5$ y $Z(8119) = Z(8115) =
2025$, el menor $n$ con al menos $2026$ ceros finales es $n = 8120$.

\textbf{23.} Base $7$: $50 = 49 + 1$, cifras (de baja a alta)
$(1, 0, 1)$. Al sumar $50 + 50$: posición $0$: $1 + 1 = 2 < 7$, sin
acarreo; posición $1$: $0 + 0 = 0$; posición $2$: $1 + 1 = 2 < 7$, sin
acarreo. Sin acarreos, luego, por Kummer, $v_7\binom{100}{50} = 0$:
$7 \nmid \binom{100}{50}$. Legendre coincide:
$v_7(100!) = \lfloor 100/7 \rfloor + \lfloor 100/49 \rfloor = 14 + 2 =
16$ y $v_7(50!) = 7 + 1 = 8$, luego
$v_7\binom{100}{50} = 16 - 2\times8 = 0$.

\textbf{24.} (i) La factorización única sostiene la definición misma
de $v_p$ y su aditividad y, por tanto, la fórmula de Legendre y toda
conclusión de \omterm{def:b1:arith:divides}{divisibilidad} (\cref{prop:b1:arith:valuation}). (ii) La
división euclídea produjo la identidad de truncamiento de la
pregunta 2 y la separación $m = p^km_1 + m_0$ que aísla el acarreo
(pregunta 12). (iii) Recuentos: el número de múltiplos de $m$
(pregunta 4), el producto de elecciones de cifras (pregunta 18) y la
cota de la suma de una fila $4^n \leq (2n+1)\binom{2n}n$
(pregunta 21) son todos argumentos al estilo
del~\cref{ch:b1:counting}.

\textbf{25.} Legendre convierte «qué potencia de $p$ \omterm{def:b1:arith:divides}{divide} a $n!$» en
aritmética de cifras en base $p$; Kummer comprime la respuesta para
los \omterm{def:b1:counting:objects}{coeficientes binomiales} en los acarreos de una sola suma --- la
\omterm{def:b1:arith:divides}{divisibilidad}, en apariencia una propiedad global de números
enormes, se lee localmente, cifra a cifra. La pregunta 16 es el
paradigma: cuatro acarreos, calculados a mano, determinan la potencia
exacta de $3$ en un número de cientos de cifras. Y la pregunta 21
muestra el mismo círculo de ideas rozando aguas profundas: una cota
inferior exponencial para $\operatorname{lcm}(1, \dots, 2n)$ es un
primer paso, enteramente elemental, hacia el teorema de los
\omterm{def:b1:arith:prime}{números primos}, cuya demostración queda muy lejos de este volumen.
Toda la caja de herramientas --- división, mcd, valuaciones --- se
repite para los polinomios en el~\cref{ch:b1:poly}, donde el análogo
del desarrollo en cifras es el desarrollo en potencias de $(X - a)$.
\end{solution}
